Activation steering---adding learned direction vectors to model hidden states---is
increasingly used to control language model behavior at inference time. Standard
practice applies steering vectors at every token position, but the temporal
dynamics of steering influence remain poorly understood: if a vector is applied
briefly and then removed, how quickly does its effect on hidden states dissipate?
We present the first direct empirical characterization of post-removal steering
vector decay. Using \gptxl and contrastive activation addition across three
behavioral dimensions (truthfulness, sentiment, formality), we measure
token-by-token cosine similarity and KL divergence after steering removal.
Our central finding is that representational decay is \emph{instantaneous}:
within a single generation step, hidden states at the steered layer revert
to values indistinguishable from an unsteered model ($\Delta$~cosine similarity
drops from ${\sim}0.019$ to ${\sim}0.000$). This result holds across all layers,
multiplier strengths, and steering durations tested. Crucially, autoregressive
and teacher-forced conditions produce \emph{identical} hidden states (Cohen's
$d = 0.000$), demonstrating that the model retains no hidden state memory of the
steering intervention. Behavioral effects persist solely through the tokens
generated under steering, which remain in the context and influence subsequent
generation. These findings establish that continuous vector application is
necessary for sustained representational steering and that the autoregressive
token sequence, not hidden state momentum, is the primary carrier of steering
influence.
